%%%%%%%%%%%%%%%%%%%%%%%%%%%%%%%%%%%%%%%%%%%%%%%%%%%%%%%%%%%%%%%%%%%%%%%%%%%%%%%%%%%
%% Roteiro de falas da apresentação do TCC                                       %%
%% Autor: Gabriel Rufino Montenegro                                              %%
%% Pareado com a apresentação Template_Beamer_UFC/document.pdf                   %%
%% Duração-alvo: ~22 min (janela 20-25 min)                                      %%
%%%%%%%%%%%%%%%%%%%%%%%%%%%%%%%%%%%%%%%%%%%%%%%%%%%%%%%%%%%%%%%%%%%%%%%%%%%%%%%%%%%
\documentclass[11pt,a4paper]{article}

\usepackage[utf8]{inputenc}
\usepackage[T1]{fontenc}
\usepackage[brazil]{babel}
\usepackage[margin=2.3cm]{geometry}
\usepackage{xcolor}
\usepackage{amsmath}
\usepackage{enumitem}
\usepackage{parskip}
\usepackage{titlesec}
\usepackage{fancyhdr}
\usepackage{hyperref}

\definecolor{ufcblue}{RGB}{0,61,124}
\definecolor{chipbg}{RGB}{224,233,243}
\definecolor{cue}{RGB}{120,120,120}

\hypersetup{colorlinks=true, linkcolor=ufcblue, urlcolor=ufcblue}

\setlength{\parskip}{0.55em}
\setlength{\parindent}{0pt}

% Cabeçalho de cada slide: número, título, faixa de tempo
\newcommand{\slide}[3]{%
  \bigskip
  \noindent\colorbox{chipbg}{\parbox{\dimexpr\linewidth-2\fboxsep\relax}{%
    \textbf{\textcolor{ufcblue}{#1}}\,\textbf{#2}\hfill\textbf{\textcolor{ufcblue}{#3}}}}\par
  \nopagebreak\medskip
}
% Indicação cênica (gesto, troca de slide, ênfase) - texto em cinza itálico
\newcommand{\cue}[1]{{\color{cue}\itshape[#1]}}

\pagestyle{fancy}
\fancyhf{}
\rhead{\small\textcolor{ufcblue}{Roteiro de falas --- TCC}}
\lhead{\small\textcolor{ufcblue}{Gabriel Rufino Montenegro}}
\rfoot{\small\thepage}
\renewcommand{\headrulewidth}{0.4pt}

\begin{document}

\begin{center}
    {\LARGE\bfseries\textcolor{ufcblue}{Roteiro de Falas da Apresentação}}\\[0.3em]
    {\large Análise comparativa de formulações de Fluxo de Potência Ótimo com ações de controle em JuMP/Julia}\\[0.2em]
    {\normalsize Gabriel Rufino Montenegro \quad$\bullet$\quad Orientador: Prof.\ Lucas Silveira Melo \quad$\bullet$\quad EE/UFC}
\end{center}

\rule{\linewidth}{0.6pt}

\textbf{Como usar este roteiro.} Cada bloco abaixo corresponde a um slide da apresentação, na mesma ordem. A faixa de tempo à direita é o \emph{tempo acumulado} ao \emph{final} daquele slide --- use um relógio/cronômetro como guia. As marcações em \cue{cinza itálico} são deixas cênicas (não devem ser lidas). O texto em preto é a fala propriamente dita; não precisa ser decorada palavra por palavra, mas o ritmo foi calibrado para fechar em torno de \textbf{22 minutos} (janela de 20 a 25). Se perceber que está atrasado, os slides de Fundamentação (Fluxo de Potência AC e Ótimo) são os mais comprimíveis.

\medskip
\textbf{Resumo do tempo:} Abertura 1{,}5 min $\to$ Fundamentação $\sim$5 min $\to$ Metodologia $\sim$4 min $\to$ Resultados $\sim$8 min $\to$ Conclusões e encerramento $\sim$3 min.

%% =====================================================================
\section*{\textcolor{ufcblue}{Bloco 1 --- Abertura}}

\slide{Slide 1.}{Capa}{0:00 -- 0:45}
Bom dia a todos. Cumprimento os membros da banca --- professor Raimundo Furtado Sampaio, professor Rafael Castro de Andrade e o engenheiro Mário Sérgio, da Intertechne ---, o meu orientador, professor Lucas Silveira Melo, e os demais presentes. Meu nome é Gabriel Rufino Montenegro e vou apresentar meu Trabalho de Conclusão de Curso, intitulado ``Análise comparativa de formulações de Fluxo de Potência Ótimo com ações de controle em JuMP/Julia''. Em uma frase: eu desenvolvi, desde a concepção, uma sequência de seis formulações de fluxo de potência que incorporam, uma a uma, as ações de controle que o operador brasileiro usa na prática, e as testei desde redes de três barras até um cenário completo do Sistema Interligado Nacional.

\slide{Slide 2.}{Sumário}{0:45 -- 1:15}
A apresentação tem quatro blocos. Primeiro, a motivação e os objetivos --- por que esse problema importa. Depois, a fundamentação teórica: o que é fluxo de potência, fluxo de potência ótimo e quais são as ações de controle. Em seguida, a metodologia, com a arquitetura de software em Julia e a descrição das seis formulações. E, por fim, os resultados sobre o conjunto de casos de teste e as conclusões. \cue{avançar para o próximo slide}

%% =====================================================================
\section*{\textcolor{ufcblue}{Bloco 2 --- Motivação e Objetivos ($\sim$2:45 de bloco)}}

\slide{Slide 3.}{Contexto: operação do SIN}{1:15 -- 2:45}
A operação de um sistema elétrico de grande porte, como o nosso Sistema Interligado Nacional, não se resume a resolver as equações de fluxo de potência. Ela depende de um conjunto de \textbf{ações de controle} que mantêm o sistema dentro de limites seguros. São cinco que me interessam aqui: o \textbf{QLIM}, que são os limites de geração de potência reativa dos geradores; o \textbf{VLIM}, os limites de magnitude de tensão nas barras; o \textbf{CSCA}, o chaveamento automático de bancos \textit{shunt}, capacitores e reatores; o \textbf{CTAP}, o ajuste automático dos \textit{taps} dos transformadores; e o que eu chamo de \textbf{DERA}, um esquema hierárquico de corte e alívio de carga em emergência. No Brasil, essas ações estão codificadas no programa ANAREDE, do CEPEL, e nos arquivos de dados no formato PWF. O ANAREDE as resolve por meio de heurísticas e laços iterativos próprios. Sem representar esses controles, um ponto de operação calculado em computador pode parecer válido e, ainda assim, estar longe do que o sistema realmente faria --- por exemplo, ignorando que um gerador já saturou seu suporte de reativo, ou que um banco de capacitores precisaria ter sido chaveado. A pergunta que move o trabalho é: \emph{como reproduzir esses controles em uma ferramenta moderna, aberta e transparente?}

\slide{Slide 4.}{Lacuna metodológica}{2:45 -- 4:00}
E aqui aparece a lacuna. Do lado acadêmico, a referência internacional é o pacote \textit{PowerModels.jl}, escrito em Julia. Ele é excelente: oferece o fluxo de potência ótimo AC e suas relaxações. Mas ele \emph{não} tem suporte nativo às ações de controle do ANAREDE. Havia uma proposta de preencher essa lacuna, o pacote \textit{ControlPowerFlow.jl}; porém, ao longo do trabalho, eu confirmei, em contato com os desenvolvedores dos pacotes PWF.jl e ControlPowerFlow, que a versão pública ficou incompleta após uma parceria com uma empresa. Ou seja, não dava para usá-lo como referência de comparação. Diante disso, minha estratégia foi desenvolver as formulações \textbf{desde a concepção}, em JuMP, adicionando os controles um a um, de modo que cada uma pudesse ser comparada contra as duas referências do PowerModels sobre o mesmo conjunto de casos de teste.

\slide{Slide 5.}{Objetivos}{4:00 -- 5:00}
O objetivo geral, portanto, foi construir, validar e comparar essa progressão incremental de formulações em JuMP/Julia, sobre o \textit{solver} Ipopt, incorporando as ações de controle do ANAREDE. Entre os objetivos específicos, destaco quatro: documentar a formulação matemática do fluxo de potência em coordenadas polares; usar o pacote \textit{PWF.jl} para ler os arquivos ANAREDE com as estruturas de controle; implementar duas formulações de referência com o PowerModels e quatro formulações próprias em JuMP; e, por fim, submeter as seis a um conjunto padronizado de vinte e sete casos de teste, do três barras ao SIN.

%% =====================================================================
\section*{\textcolor{ufcblue}{Bloco 3 --- Fundamentação Teórica ($\sim$4:30 de bloco)}}

\slide{Slide 6.}{Fluxo de Potência AC}{5:00 -- 6:15}
Rapidamente, a base teórica. O fluxo de potência AC é o conjunto de equações de balanço de potência ativa e reativa em cada barra: o que entra de geração, menos o que sai de carga, tem que fechar com o que circula pelos ramos. As barras são classificadas em três tipos: a \textit{slack}, que é a referência de ângulo; as PV, com tensão e potência ativa especificadas; e as PQ, com as potências dadas e tensão e ângulo como incógnitas. O ponto importante para o restante da apresentação é que esse sistema é \textbf{não linear e não convexo} --- ele é tradicionalmente resolvido por Newton-Raphson, e a não convexidade é justamente o que torna o problema difícil em larga escala.

\slide{Slide 7.}{Fluxo de Potência Ótimo e relaxações}{6:15 -- 7:45}
Quando a esse sistema eu adiciono uma função objetivo --- tipicamente o custo de geração --- e o submeto a restrições, eu tenho o fluxo de potência \emph{ótimo}. Lendo a equação do slide: eu minimizo o custo total de geração, \emph{sujeito a} --- é o que significa o ``s.a.'' --- que o balanço de potência feche e que os limites operacionais sejam respeitados; esses limites incluem as faixas de tensão, os limites de geração e o limite térmico das linhas. Uma observação que retomo mais à frente: como no meu caso não há dados reais de custo das fontes, minimizar custo com coeficientes iguais equivale a minimizar a geração total e, pela conservação de energia com carga fixa, a minimizar as perdas. Como ele herda a não convexidade, surgiram na literatura várias relaxações. Há uma distinção conceitual que vale frisar: a aproximação \textbf{DC} lineariza as equações --- ela é rápida, mas \emph{não} dá garantia sobre o ótimo verdadeiro; já as relaxações \textbf{SOC} e \textbf{SDP} são, por construção, garantidas: fornecem uma cota inferior \emph{válida} do custo real. A tabela mostra isso num caso de cinco barras: o DC fica \emph{abaixo} do AC porque ignora perdas e impõe condições irreais; e o SOC fica ainda mais abaixo, com um \textit{gap} de cerca de dezoito por cento, o que é esperado em redes malhadas. Esse estudo é apenas didático --- a progressão principal do trabalho mantém a formulação AC exata.

\slide{Slide 8.}{Ações de controle do ANAREDE}{7:45 -- 9:30}
Este é o slide conceitualmente central. Eu represento cada controle assim: o \textbf{QLIM e o VLIM} reproduzem a lógica de troca de tipo de barra, PV para PQ, mas em vez de chavear de forma descontínua --- o que atrapalha um \textit{solver} suave como o Ipopt ---, eu uso \textbf{variáveis de folga penalizadas} na função objetivo. O \textbf{CSCA} transforma a susceptância do banco \textit{shunt} numa variável contínua de decisão, dentro de seus limites. O \textbf{CTAP} faz o mesmo com a razão de \textit{tap} do transformador, e reaproveita as folgas de tensão já existentes, sem criar folga nova. E o \textbf{DERA} é um corte hierárquico de carga: o \textit{solver} pode cortar parte da carga ativa de cada barra, mas com penalidade crescente conforme o nível de tensão --- ele corta primeiro onde é mais fácil religar e protege a extra-alta tensão. O fio condutor é que \emph{todos} os controles entram como \textit{soft-constraint}, ou seja, restrições flexíveis penalizadas, evitando os laços descontínuos do ANAREDE. E cada folga tem significado físico claro: a folga de tensão do QLIM, por exemplo, corresponde ao gerador saturando o regulador automático de tensão e travando no limite de capabilidade --- a tensão da barra só se afasta do \textit{setpoint} quando isso acontece, que é exatamente a conversão de barra PV para PQ que o ANAREDE executa no seu laço de troca de tipo. Já a folga de reativo do VLIM permite que a demanda reativa da barra de carga se ajuste minimamente para manter a tensão dentro dos limites.

%% =====================================================================
\section*{\textcolor{ufcblue}{Bloco 4 --- Metodologia ($\sim$3:45 de bloco)}}

\slide{Slide 9.}{Fluxo de execução}{9:30 -- 10:45}
Do ponto de vista de software, tudo foi feito em Julia, encadeando quatro pacotes: o \textit{PWF.jl} para ler o arquivo ANAREDE, o \textit{PowerModels.jl} como referência, o \textit{JuMP.jl} para escrever os modelos e o \textit{Ipopt} como \textit{solver} não linear. O fluxo de execução tem seis etapas, sempre as mesmas: leitura do PWF já com os dados de controle; uma adequação da topologia que mantém só a maior componente conexa e uma única barra de referência; construção da estrutura de referência; montagem do modelo no JuMP; resolução; e exportação dos resultados em CSV. Um detalhe de modelagem: os elos de corrente contínua, os HVDC, têm suas injeções \textbf{fixadas} como condição de contorno --- elas já vêm do fluxo base, e liberá-las só aumentaria o problema sem trazer informação.

\slide{Slide 10.}{A progressão F0 $\to$ F5}{10:45 -- 12:00}
Esta tabela resume as seis formulações. A \textbf{F0} é o fluxo de potência ótimo de referência, do PowerModels; a \textbf{F1} é o fluxo de potência clássico, também do PowerModels. Daí em diante começam as minhas formulações em JuMP: a \textbf{F2} acrescenta QLIM e VLIM; a \textbf{F3} acrescenta o CSCA; a \textbf{F4}, o CTAP; e a \textbf{F5}, o DERA. O ponto-chave é que a progressão é \textbf{estritamente aditiva}: cada formulação parte da anterior e acrescenta um único bloco de variáveis e restrições. Isso tem uma consequência metodológica forte --- se duas formulações consecutivas divergem num resultado, essa divergência é atribuível \emph{exclusivamente} ao novo controle introduzido. É o que me permite isolar o efeito de cada controle.

\slide{Slide 11.}{Como funcionam as folgas penalizadas}{12:00 -- 13:15}
Para deixar concreto, veja a F2. As variáveis primárias são as de sempre --- tensão, ângulo, potências ---, com os limites lidos do PWF. A novidade são as folgas: uma folga de tensão em cada barra de geração --- o mecanismo do QLIM --- e uma folga de reativo em cada carga --- o mecanismo do VLIM. A restrição de \textit{setpoint} diz que a tensão da barra é o valor de referência mais a folga; e a função objetivo minimiza a soma dos quadrados dessas folgas, com um peso alto, de um milhão. A intuição é simples e é o coração do trabalho: \emph{se} o problema é factível respeitando os \textit{setpoints}, a folga fica zero e a solução respeita tudo. \emph{Se não é}, em vez de o \textit{solver} simplesmente declarar ``infactível'', como faz o fluxo clássico, ele encontra a solução que \textbf{viola minimamente} os controles. É essa robustez que vai aparecer nos resultados.

%% =====================================================================
\section*{\textcolor{ufcblue}{Bloco 5 --- Resultados ($\sim$8:00 de bloco)}}

\slide{Slide 12.}{Conjunto de casos de teste --- 27 casos}{13:15 -- 14:15}
O conjunto tem vinte e sete casos de teste, escolhidos para varrer dois eixos: o porte da rede e a presença de cada controle. Há redes pedagógicas de três a nove barras, com variantes que ativam cada seção do PWF; redes de médio porte de trezentas e quinhentas barras; dois recortes reduzidos do Nordeste do SIN, de mil e de dois mil e seiscentas barras, derivados de um caso real do ciclo de planejamento do ONS; o cenário integral do SIN, com mais de treze mil barras; e alguns casos sintéticos de estresse. O critério para dizer que duas formulações chegaram à ``mesma'' solução é rígido: diferença menor que dez à menos quatro, simultaneamente, em tensão, ângulo, potências e fluxos.

\slide{Slide 13.}{Convergência global}{14:15 -- 15:45}
E este é o primeiro resultado central. Olhando a faixa de convergência no conjunto completo de casos: minhas formulações desenvolvidas manualmente, da F2 à F5, convergem em \textbf{vinte e três} dos vinte e sete casos. As referências convergem em menos: dezoito para a F1 e quinze para a F0. A tabela colorida mostra uma amostra --- repare nas linhas em amarelo, como o \texttt{4busfrank\_vlim}, o \texttt{5busfrank\_csca} e o \texttt{caso\_red2}: ali as referências declaram infactibilidade, e as formulações controladas convergem. No total, há \textbf{nove} casos em que pelo menos uma referência falha e pelo menos uma das minhas formulações encontra solução. A razão é exatamente aquela das folgas: onde o limite operacional rígido inviabiliza o problema, a folga penalizada recupera um ponto de operação utilizável. E quero ser preciso sobre o que isso significa: não é que minhas formulações sejam ``melhores'' por convergir mais; é que elas respondem a uma pergunta diferente --- em vez de apenas dizer ``infactível'', elas dizem \emph{o quanto} e \emph{onde} o sistema precisaria violar os controles para conseguir operar.

\slide{Slide 14.}{Análise barra a barra --- \textit{clusters}}{15:45 -- 17:00}
Convergir é uma coisa; \emph{para onde} cada formulação converge é outra. Aqui eu agrupo as soluções em \textit{clusters} de equivalência. Veja a rede de três barras: aparecem quatro soluções distintas. A \textbf{F0}, o ótimo do FPO --- que, na ausência de dados reais de custo das fontes, equivale na prática a minimizar as perdas ativas ---, empurra a tensão para o limite de 1,1 por unidade. A \textbf{F1} dá o ponto operacional ``natural'' da rede. A \textbf{F2} fica perto da F1, com uma folguinha residual. E as \textbf{F3 a F5}, que têm o CSCA ativo, deslocam o ponto para tensões consistentemente mais baixas, porque o controle do banco \textit{shunt} ajusta a susceptância para perseguir o \textit{setpoint}. Ou seja, cada \textit{cluster} é uma assinatura computacional de um controle --- o que serve, inclusive, para detectar regressões em implementações futuras.

\slide{Slide 15.}{Destaque: o \texttt{caso\_red2}}{17:00 -- 18:20}
Esse caso é o exemplo mais nítido de todo o conjunto de casos de teste, e eu gosto de detê-lo aqui. É um recorte real do Nordeste, com mil e trinta e três barras. O fluxo de potência clássico, a F1, declara infactibilidade --- mas repare \emph{como}: ele entrega uma ``solução'' com tensão mínima de \textbf{menos zero vírgula quatro} por unidade e máxima de quase dois por unidade. Isso é fisicamente impossível. Já as formulações controladas, da F2 à F5, convergem para tensões dentro de uma faixa realista. E mais: a F4, que tem CSCA mais CTAP, reduz a função objetivo a cerca de dez à menos dois --- algo como cinco mil vezes menor que a F2. Na prática, a combinação dos dois controles encontra um ponto de operação praticamente compatível com os \textit{setpoints} do arquivo, quase sem violação. É a demonstração concreta de para que servem essas formulações.

\slide{Slide 16.}{O caso SIN integral}{18:20 -- 19:35}
Agora, o caso de maior interesse prático: o SIN integral, com treze mil trezentas e trinta e oito barras. Aqui eu preciso ser honesto sobre um resultado \emph{negativo}: \textbf{nenhuma} das seis formulações converge dentro do orçamento de três mil iterações e do tempo limite. A F0, a F2 e a F4 estouram o tempo; a F1, a F3 e a F5 terminam em infactibilidade local. Mas há uma leitura importante junto: o sistema \emph{não quebra} --- toda a infraestrutura, leitura do PWF, construção da referência, montagem do modelo e chamada ao Ipopt, funciona corretamente em escala SIN. Não há erro de \textit{parser} nem \textit{crash}. E um ponto importante: minhas formulações \emph{não} partem de uma inicialização trivial --- elas já são inicializadas com o ponto de operação convergido pelo próprio ANAREDE, lido do arquivo PWF. Mesmo assim não fecham; e a F1, a referência do PowerModels, partindo do mesmo ponto, também termina infactível. Ou seja, o gargalo não está na inicialização: está na fidelidade do modelo reconstruído --- elos HVDC, intercâmbios de área, controles remotos, múltiplas barras de referência --- e no condicionamento numérico do problema nessa escala. O obstáculo é \textbf{numérico}, não estrutural.

\slide{Slide 17.}{Verificação auxiliar FDC}{19:35 -- 20:50}
E para comprovar essa hipótese eu fiz uma verificação adicional, que chamo de FDC: resolver o \emph{mesmo} arquivo do SIN, mas na aproximação DC, linear. O resultado é o oposto: ele converge em poucas iterações, com status resolvido. A tabela traz os números agregados --- treze mil barras retidas, cerca de setenta e um gigawatts de geração, perdas nulas por definição. Isso \textbf{confirma} que o obstáculo do AC é numérico, não estrutural: a mesma rede, no mesmo fluxo de execução, fecha quando linearizada. Há um porém honesto: os ângulos da solução DC chegam a cento e setenta graus, o que viola as próprias hipóteses da linearização --- então essa solução não é fisicamente realizável e, com tensão um por unidade em todas as barras, ela seria um ponto de partida \emph{pior} que o ponto convergido do ANAREDE que eu já uso. Por isso o valor da FDC é \emph{diagnóstico}: ela confirma que existe solução para a topologia e que o obstáculo do AC é numérico --- o caminho de trabalho futuro é reconciliar o modelo com o do ANAREDE e melhorar o condicionamento, não trocar o ponto de partida.

%% =====================================================================
\section*{\textcolor{ufcblue}{Bloco 6 --- Conclusões e encerramento ($\sim$3:00 de bloco)}}

\slide{Slide 18.}{Síntese das contribuições}{20:50 -- 22:20}
Em síntese, três contribuições. A primeira: as formulações controladas ampliam a faixa de convergência --- vinte e três de vinte e sete casos, recuperando soluções factíveis onde as referências falham. A segunda: os \textit{clusters} de equivalência funcionam como uma ``impressão digital'' computacional de cada controle, reproduzindo o comportamento esperado do QLIM, do CSCA, do CTAP e do DERA. E a terceira, que é uma contribuição negativa mas relevante: eu delimitei com clareza a fronteira atual --- nenhuma formulação converge no SIN, mesmo partindo do ponto já convergido pelo ANAREDE, o que mostra exatamente onde a abordagem \textit{soft-constraint} pura precisa de reconciliação de modelo e de melhor condicionamento numérico. Em conjunto, essas três contribuições mostram que a abordagem é robusta nas escalas pedagógica e regional, e mapeiam, com precisão, o que ainda falta para alcançar a escala nacional.

\slide{Slide 19.}{Limitações e trabalhos futuros}{22:20 -- 23:35}
As principais limitações: o CSCA e o CTAP são tratados como variáveis contínuas, enquanto a operação real é discreta; o DERA é uma versão contínua e suave de um mecanismo que, na realidade, é discreto e emergencial; e os elos HVDC entram com fluxos fixados. Os trabalhos futuros decorrem diretamente disso: primeiro, viabilizar o SIN atacando a causa efetiva --- uma verificação de resíduos alimentando o modelo com a solução do PWF, uma auditoria da tradução PWF para JuMP e o escalonamento progressivo das penalidades por homotopia; segundo, discretizar o CSCA e o CTAP, transformando o problema num MINLP, com os \textit{solvers} que o ecossistema Julia já oferece; e, terceiro, acoplar as minhas formulações ao \textit{ControlPowerFlow.jl} como ferramenta de validação.

\slide{Slide 20.}{Encerramento}{23:35 -- 24:05}
Com isso eu encerro. O trabalho mostrou que é possível reproduzir, de forma transparente e reprodutível, as ações de controle do ANAREDE em uma ferramenta aberta, e mapeou tanto os ganhos quanto a fronteira numérica atual da abordagem. Agradeço novamente à banca pela atenção e ao professor Lucas pela orientação, e fico à disposição para as perguntas. Muito obrigado.

%% =====================================================================
\bigskip
\rule{\linewidth}{0.6pt}
\section*{\textcolor{ufcblue}{Apêndice --- Perguntas prováveis da banca (preparação)}}

\textbf{1. Por que penalidade quadrática e por que $\rho = 10^6$?} O quadrado mantém a função objetivo suave e diferenciável, ideal para o Ipopt; o valor $10^6$ é alto o bastante para que, havendo solução sem violação, as folgas saturem em zero, mas não tão alto a ponto de degradar o condicionamento numérico. Na F5 eu reduzo para $10^5$ justamente para que o termo de corte de carga se torne competitivo.

\textbf{2. Tratar CSCA/CTAP como contínuos não descaracteriza o controle?} Captura corretamente a \emph{direção} e a \emph{magnitude} da compensação que o operador faria; perde-se só a granularidade discreta dos passos. A discretização via MINLP é o passo seguinte, e está nos trabalhos futuros.

\textbf{3. Por que o SIN não converge se os recortes convergem?} Não é a inicialização: todas as formulações já partem do ponto convergido pelo ANAREDE, lido do PWF --- e mesmo a F1, do PowerModels, partindo desse ponto, termina infactível. A diferença está no que só o caso integral tem: elos HVDC, intercâmbios de área, controles remotos e múltiplas barras de referência, cuja tradução PWF$\to$JuMP precisa ser auditada, além do condicionamento do NLP com penalidade $10^6$ nessa escala. A FDC confirma que o obstáculo é numérico, não estrutural --- por isso as frentes são verificação de resíduos, auditoria da tradução e homotopia das penalidades.

\textbf{4. Qual a diferença essencial para o ANAREDE?} O ANAREDE resolve por laços iterativos externos e chaveamentos discretos; eu resolvo tudo numa única otimização não linear com folgas. Os resultados convergem para o mesmo tipo de ponto de operação, mas com uma formulação explícita e auditável.

\textbf{5. O que é exatamente o DERA neste trabalho?} É um rótulo herdado do nome do \textit{script}. Aqui ele designa um esquema hierárquico de corte de carga, análogo aos Esquemas Regionais de Alívio de Carga do ONS, em que o \textit{solver} corta carga preferencialmente nas barras de menor nível de tensão.

\textbf{6. Por que Julia e não Python/MATLAB?} Pela combinação de sintaxe de alto nível com desempenho de linguagem compilada, e porque o JuMP e o PowerModels.jl, referências em otimização de sistemas de potência, são nativos do ecossistema Julia.

\textbf{7. Por que a F0 minimiza perdas se o objetivo é custo?} Pela conservação de energia no sistema: geração total $=$ carga total $+$ perdas. A carga é dado fixo de entrada, então geração e perdas diferem apenas por uma constante --- minimizar uma minimiza a outra. Como eu não tenho dados reais de custo das fontes e atribuo coeficientes iguais a todos os geradores, minimizar o custo total vira minimizar a geração total, que por sua vez é minimizar as perdas. Por isso o resultado da F0 deve ser lido como um ponto de mínimas perdas, não de mínimo custo econômico --- e é essa a ressalva metodológica que registro na dissertação.

\textbf{8. Na equação de balanço, por que o termo \textit{shunt} é $V_i^2$ e não $V_i V_j$?} Porque são elementos diferentes. O $g^{sh}_i V_i^2$ é um \textit{shunt} ligado entre a barra $i$ e o \emph{terra} --- como o outro lado é o terra, com tensão zero, só entra $V_i$, e o subíndice é $i$ porque o elemento está preso a uma barra só. O acoplamento $V_i V_j$ que se espera existe, sim, mas dentro do fluxo de ramo $p_{ij}$: ao abrir esse termo aparecem tanto um termo próprio $V_i^2$ quanto o termo mútuo $V_i V_j$, agora com as admitâncias de subíndice $ij$, porque aí o elemento é a linha entre as duas barras.

\medskip
\textit{\textcolor{cue}{--- Perguntas específicas de otimização ---}}

\textbf{9. Qual o método de otimização utilizado?} O \textit{solver} é o Ipopt, que implementa um método de \textbf{pontos interiores primal-dual} com busca em linha do tipo \textit{filter} (Wächter e Biegler, 2006). A ideia central: as restrições de desigualdade (os limites) são incorporadas à função objetivo como termos de \emph{barreira logarítmica} $-\mu \sum \ln(\text{folga})$, que impedem as iterações de tocar as fronteiras; o \textit{solver} resolve uma sequência de subproblemas, reduzindo o parâmetro de barreira $\mu \to 0$, e a cada passo aplica o método de Newton sobre o sistema de condições de otimalidade (KKT) perturbadas, o que envolve resolver um sistema linear esparso. É por exigir derivadas de primeira e segunda ordem que o método precisa de funções \emph{suaves} --- contínuas e diferenciáveis ---, e é isso que motiva toda a estratégia de folgas em vez de lógicas discretas.

\textbf{10. O problema é convexo? Como garantir o ótimo global?} Não --- o FPO AC é \textbf{não convexo}, por causa dos termos $V_i V_j \cos(\theta_i-\theta_j)$. O Ipopt, por ser um método local, encontra apenas um \emph{ponto KKT local}, sem garantia de ótimo global. Eu mitigo isso inicializando as formulações no ponto de operação já convergido pelo ANAREDE, que é fisicamente sensato. Garantia global exigiria \textit{branch-and-bound} espacial ou o uso de relaxações convexas (SOC, SDP) para fornecer uma cota inferior e limitar o \textit{gap} de otimalidade --- foi justamente o que o estudo do \texttt{case5.m} ilustrou.

\textbf{11. Como o Ipopt trata os limites de desigualdade?} Pelo termo de barreira: em vez de testar quais restrições estão ativas (como um método de conjunto ativo ou o \textit{simplex}), ele mantém as iterações estritamente no \emph{interior} da região factível e deixa que a redução de $\mu$ as empurre suavemente contra os limites que forem ativos na solução. Por isso o nome ``pontos interiores''.

\textbf{12. O que significam os \textit{status} LOCALLY\_SOLVED e LOCALLY\_INFEASIBLE?} \textbf{LOCALLY\_SOLVED} quer dizer que o Ipopt encontrou um ponto que satisfaz as condições KKT de primeira ordem dentro da tolerância --- um ótimo \emph{local}. \textbf{LOCALLY\_INFEASIBLE} quer dizer que ele convergiu para um mínimo local da \emph{violação} das restrições sem atingir factibilidade; é importante frisar que isso \emph{não} prova que o problema é globalmente infactível --- apenas que, a partir daquele ponto de partida, o \textit{solver} não encontrou uma região factível.

\textbf{13. Por que penalizar (\textit{soft-constraint}) em vez de impor as restrições diretamente?} Porque impor a tensão igual ao \textit{setpoint} como igualdade rígida, \emph{junto} com os limites rígidos de reativo, pode ser um sistema incompatível --- sobredeterminado --- e o \textit{solver} devolve simplesmente INFEASIBLE. A penalidade relaxa o \textit{setpoint} de forma controlada, garantindo um problema sempre factível que se degrada suavemente: em vez de um modelo que ``quebra'', tenho um que entrega a melhor solução possível com violação mínima. É essa a diferença que aparece nos casos em que só as minhas formulações convergem.

\textbf{14. A penalidade quadrática é ``exata''? As folgas realmente chegam a zero?} No sentido técnico, não: penalidade \emph{exata} exige uma norma não suave (tipo $L_1$), que o Ipopt não digere bem. Com a penalidade quadrática e $\rho$ finito, as folgas tendem a zero apenas no limite $\rho \to \infty$; com $\rho = 10^6$ sobra um resíduo numericamente desprezível. Aumentar $\rho$ reduz o resíduo mas piora o condicionamento --- o escalonamento progressivo de $\rho$ por homotopia é o caminho registrado nos trabalhos futuros.

\textbf{15. Se discretizasse o CSCA e o CTAP, que classe de problema e que método?} Viraria um \textbf{MINLP} --- programação não linear inteira mista, com variáveis inteiras (o passo do banco, a posição do \textit{tap}) convivendo com as equações não lineares da rede. A resolução típica é por \textit{branch-and-bound}, que resolve uma sequência de relaxações contínuas --- e as minhas F3 e F4 são exatamente essas relaxações. O ecossistema Julia já oferece \textit{solvers} acoplados ao JuMP para isso: Juniper, Bonmin, Alpine ou SCIP. O custo computacional, porém, cresce muito.

\textbf{16. Se a FDC (DC) converge sobre o SIN, por que a solução não presta --- e por que não a uso como \textit{warm-start}?} Porque a aproximação DC só é válida sob hipóteses que essa própria solução viola. O DC lineariza o fluxo ativo trocando $\sin(\theta_i-\theta_j)$ por $(\theta_i-\theta_j)$, o que só vale para \emph{diferenças angulares pequenas}. A troca é excelente a 10$^\circ$ (erro $\sim$0\%), sofrível a 30$^\circ$ ($\sim$5\%) e absurda a 170$^\circ$: o transporte real, $\sin 170^\circ = 0{,}17$, é pequeno, mas o modelo linear usa 2{,}97 rad e superestima o fluxo em cerca de 17 vezes. É por isso que aparece um fluxo de ramo de 99{,}5\,p.u., quase 10\,GW numa linha --- um valor que nenhuma linha real carregaria, pois passando de $\sim$90$^\circ$ por linha o sistema síncrono já perde estabilidade estática. Então a solução DC é válida para o \emph{modelo linearizado}, mas não corresponde a um ponto de operação fisicamente realizável da rede AC. Como ela ainda tem tensão 1{,}0 em todas as barras e ângulos de $\pm 177^\circ$, é um ponto de partida \emph{pior} que o ponto já convergido do ANAREDE que eu utilizo --- por isso o papel da FDC é \textbf{diagnóstico} (confirma que existe solução para a topologia, logo o obstáculo do AC é numérico, não estrutural), e não de inicialização. \textit{\textcolor{cue}{(Se cutucarem: os $\pm 170^\circ$ são ângulos de barra em relação à \textit{slack}, que acumulam ao longo de muitas linhas em série no SIN; a prova de que \emph{alguma} linha está fora da faixa de pequeno ângulo é o próprio fluxo de 99{,}5\,p.u., pois no DC $p_{ij}=\Delta\theta_{ij}/x_{ij}$.)}}

\textbf{17. O que quer dizer ``o obstáculo do SIN é numérico, não estrutural''?} São duas causas distintas para uma mesma falha de convergência. \emph{Estrutural} significaria que a solução \textbf{não existe} --- a rede seria genuinamente infactível, e aí nada adiantaria (nem inicialização, nem ajuste de \textit{solver}). \emph{Numérico} significa que a solução \textbf{existe, mas o algoritmo não chega até ela} por dificuldade de cálculo: mau condicionamento, escala, penalidades desequilibradas. A analogia é um labirinto: estrutural é um labirinto \emph{sem saída}; numérico é um labirinto \emph{com saída}, mas você está de olhos vendados e na neblina, e o tempo acaba antes. No meu caso, três evidências sustentam que é numérico: (i) o \textbf{ANAREDE já achou um ponto de operação real} dessa rede --- então uma solução física existe; (ii) a \textbf{FDC (DC) converge} sobre a mesma topologia, logo há um modelo solúvel daquela rede; e (iii) mesmo \textbf{partindo do ponto bom}, todas as seis formulações e até a F1 do PowerModels falham --- o que aponta para a fidelidade da tradução PWF$\to$JuMP (elos HVDC, intercâmbios, controles remotos, múltiplas referências) e para o condicionamento do NLP com $\rho = 10^6$ em treze mil barras, não para a ausência de solução. Isso muda o desfecho do trabalho: em vez de um beco sem saída, o caminho à frente é engenharia numérica --- verificação de resíduos, auditoria da tradução e homotopia de $\rho$ ---, não uma teoria nova.

\textbf{18. Se a FDC não tem perdas, por que a geração ($\approx$70{,}97\,GW) difere da carga ($\approx$69{,}57\,GW)?} Essa diferença de $\sim$1{,}4\,GW \textbf{não é perda} --- em DC as perdas são exatamente zero por definição, e o balanço nodal completo fecha na igualdade. É um efeito de \emph{contabilidade}: os elos HVDC do SIN (Itaipu, Madeira, Belo Monte) injetam e absorvem potência ativa real na rede AC, modelados como geradores fictícios de injeção fixada, mas essas injeções \emph{não entram} na coluna \texttt{P\_Geracao} do CSV que eu uso para totalizar a geração --- elas entram apenas no balanço interno do \textit{PowerModels}. Somando ainda o pequeno ajuste absorvido pela(s) barra(s) de referência, quando se contabiliza a contribuição dos elos e da \textit{slack} o balanço fecha exatamente. Ou seja, o 1{,}4\,GW mede o que trafega pelos elos HVDC, não dissipação --- é um artefato de qual coluna do CSV está sendo somada, não uma perda física.

\end{document}
